\begin{otherlanguage}{english}
\noindent\textbf{Fibrations, localizers, and the coherence of fibres}

\medskip
The pages assembled here turn a useful geometric intuition into a problem of foundations. A family of categories should behave like a fibration when its strict fibres calculate homotopy fibres and vary coherently over the base. Grothendieck first characterizes that behaviour for ordinary weak homotopy equivalences. He then asks which axioms on an abstract class \(W\) of functors permit the same constructions, and discovers that several apparently natural requirements already force the whole ordinary homotopy theory. The resulting revisions are part of the argument: the word “fibration” changes its tests as the desired scope of the theory becomes clearer.

These pages comprise Chapter XI, manuscript pp.~127--136, and Chapter XII, manuscript pp.~1--58, through the end of section~IV. These are separate manuscript sequences. The edition explicitly says that Chapter XI has no pp.~1--126; their absence is not a gap in the text. We write \(BC=|NC|\) for the classifying space of a category and \(W_\infty\) for the functors whose nerves are weak homotopy equivalences.

\medskip
\noindent\textbf{The geometric theorem and its coefficient proof.}
Chapter XI begins with a property of the actual pullback functor:
\[
 f^*:\Cat/Y\longrightarrow\Cat/X
 \quad\hbox{preserves }W_\infty.
\]
This is a universal condition on changes of base, not a condition merely on the homotopy types of the individual fibres. Theorem~1 identifies it with universal comma-category tests and with base-change and local-constancy statements for direct images of locally constant coefficients. In the terminology later used for model categories, these are sharp maps for the Thomason weak equivalences \citep[\S2]{Rezk1998Sharp}. Quillen's Theorem B explains the geometry of the comma tests: when their transition functors are weak equivalences, the strict fibre, compared with the appropriate comma category, represents the homotopy fibre \citep[\S1]{Quillen1973}.

Most of the coefficient argument is complete as written. The degree-zero comparison and the first-quadrant Leray comparison require no new finiteness hypothesis. The nonabelian passage needs more care. A component of a torsor groupoid has the automorphism group of a representative torsor; this group can be an inner form of the coefficient group, and its direct image is not generally the fixed group printed in the proposed obstruction sequence. We therefore replace that sequence by an argument with the groupoids of torsors on \(X/y\). The universal degree-zero hypothesis gives full faithfulness of their transition functors; the degree-one hypothesis gives essential surjectivity. The same two hypotheses identify their values with the torsor groupoids on strict fibres. Descent reconstructs global torsors from this locally constant diagram, and homotopy invariance of such diagrams proves the required comparison \citep[Theorem~1.3]{Cisinski2009LocallyConstant}.

Theorem~2 then detects an equivalence between two such families on one fibre in each connected component of the base. The source's local reduction for torsors becomes precise through the same descent argument. Its converse uses a contractible replacement of the base. We give an independent construction by factoring a point in Thomason's model category \citep[Theorem~4.9]{Thomason1980}, and also check the specific path-category dependency that Chapter XII recalls from Chapter VII. This separates the existence of a suitable replacement from the stronger task of constructing it explicitly with both properness and smoothness.

\medskip
\noindent\textbf{What the axioms force.}
Chapter XII starts with identities and cancellation, a theorem-A principle, and the asphericity of categories with a terminal object. Two forms of theorem A occur. W(2) is relative: equivalences on comma categories over every object of a third category imply equivalence of the total functor. W(2 bis) is absolute: aspherical comma categories imply that a functor is an equivalence. The manuscript deliberately uses the latter in several arguments. Modern “fundamental localizer” normally includes the former. We keep the distinction throughout the audit.

The minimum conjectured in these pages is now known. Cisinski identifies \(W_\infty\) as the smallest fundamental localizer \citep[Theorem~2.2.11]{Cisinski2004Minimal}, and as the smallest weak fundamental localizer with the absolute theorem-A axiom \citep[Theorem~6.1.18]{Cisinski2006}. Strong saturation and stability under filtered colimits are also established for full fundamental localizers; those results are not assertions about every class satisfying only the weaker axiom. An inference from “invertible in the localization” to “belongs to \(W\)” is justified at the point where its applicable hypothesis is available.

The fibration axioms add a rigidity that the basic localizer axioms do not possess. Their weakest perfect-fibre form, combined with nondegeneracy, makes homotopy fibres of \(W\)-maps \(W\)-aspherical. Iterating the loop construction then detects all homotopy groups. We repair the source's nondegeneracy argument without its unjustified restriction to a component: a missed component produces an empty homotopy fibre, while two identified components produce a disconnected aspherical fibre, whose diagonal in its square again has an empty homotopy fibre. Asphericity of the empty category collapses the theory. Thus the alternatives are \(W_\infty\) and the class of every functor. The second alternative cannot be discarded from the opening assertion about derivators; the terminal derivator realizes it.

The truncation examples explain why a more flexible fibration notion is needed. The class \(W_n\) retains homotopy groups only through degree \(n\), and is a genuine localizer. But the based-path family of an \((n+1)\)-sphere has an invisible total space and an invisible base while its fibre retains a nonzero \(n\)-th homotopy group, or a nontrivial component set when \(n=0\). Hence the strongest fibre axioms fail. In degree zero, the exact sequence must be read in pointed sets and with the appropriate group action: a cokernel of groups does not describe the component set in general.

\medskip
\noindent\textbf{Lift categories supply the missing structure.}
Section~III develops proper and smooth functors through categories of lifts with prescribed source or target. Those categories turn universal comma comparisons into workable tests, prove stability under base change and composition, and explain why Grothendieck (op)fibrations furnish examples. They also support the passage from fibrewise equivalence to equivalence of total categories. The later formulations by Maltsiniotis describe the corresponding exact squares, while the derivator account explains the associated base-change transformations \citep{Maltsiniotis2005Asphericite,Maltsiniotis2012Exact,Groth2013}.

The proposed weakening to endpoint equivalences over every arrow loses a necessary compatibility. We give a poset with eight objects over \([2]\). Its fibres have two contractible components, every interval pullback has both endpoint inclusions in \(W_\infty\), and transport along \(0\to1\) and \(1\to2\) preserves the component labels. Transport along \(0\to2\) exchanges them. Its total nerve is an eight-cycle. Consequently the edge transports do not compose, and the pseudofibration question raised on manuscript p.~40 has a negative answer. The complete proof is independent of the accompanying finite computation.

The successful replacement uses universal preproperness or prelissity in addition to the interval test. Comma categories then provide actual functorial transport. A split replacement converts this into a diagram of categories whose arrows lie in \(W\). For a categorically contractible base, a finite zigzag of natural transformations gives an explicit finite zigzag of fibrewise and total \(W\)-equivalences. This proves the contractible-base portion of Theorem~3 directly, without extracting a zigzag from an abstract isomorphism in a localization. Here “contractible” has the manuscript's categorical meaning; a weakly contractible nerve need not admit such a finite contraction.

\medskip
\noindent\textbf{The last proposition and the limits of the repair.}
Section~IV returns to the relations among the fibration axioms, then proposes a fibrewise gluing statement. Its final Proposition~3 is false literally as printed. Over \([1]\), take the strictly Cartesian inclusion of the split fibration \(\{o\to a\}\) into \(\{o\to a,o\to b\}\). Every fibre map is a functor between discrete categories and is perfect. Pulling the total functor back along \(o\to b\) gives endpoint fibres \(e\) and \(\varnothing\); it is not even a pseudofibration. The target structural family is not a \(W_\infty\)-fibration, precisely the extra condition mentioned in the preceding discussion and omitted from the proposition.

We prove a sufficient repair for that preceding, fibrewise-perfect situation: the structural transports lie in \(W\), the fibre maps are smooth \(W\)-fibrations, and W(4a) and W(4b) control the relevant pullback and fibre comparison. We do not claim that adding only the omitted structural condition repairs the full printed statement. Rezk's later gluing theorem has an analogous additional requirement that naturality squares be homotopy Cartesian \citep[Theorem~7.1]{Rezk1998Sharp}; passing between models does not by itself remove the need to control strict pullbacks.

The marked developments give these deductions with complete proofs and bounded accounts of the literature examined, without a claim of priority. The registers distinguish what the source proves, sketches, assumes, or conjectures from the status supplied by the commentary. They preserve edition slips in the French text and apply only corrections of the transcription. Questions not decided under their exact weaker hypotheses are identified as limitations of this audit, rather than declared open in the literature.

\medskip
\noindent\textit{Commentary and mathematical audit AI-produced by GPT-6 Astra (OpenAI, ChatGPT Pro mode), checked and integrated by Claude (Anthropic).}
\nocite{*}
\end{otherlanguage}
